\subsection{SPACES OF COMPLEX PROJECTIVE LINEAR VARIETIES}

With the notation recalled in Chapter VI, § 3, no. 5, we define similarly the spaces of projective linear varieties in a complex projective space:

DEFINITION 2. The quotient space \(\mathbf{P}_{n,p}(\mathbf{C})\) of the topological space \(\mathbf{L}_{n+1, p+1}(\mathbf{C})\) by the equivalence relation \(\Delta_{n,p}(\mathbf{C})\) is called the space of projective linear varieties of \(p \geqslant 0\) dimensions in the projective space \(\mathbf{P}_n(\mathbf{C})\) .

By the argument of Proposition 6 of Chapter VI, § 3, no. 5, we see first of all that \(\mathbf{P}_{n,p}(\mathbf{C})\) is Hausdorff. Next we prove that it is compact by replacing the subspace \(\mathrm{V}_{n+1,p+1}\) (in the proof of Proposition 7 of Chapter VI, § 3, no. 5) by the subspace \(\mathrm{W}_{n+1,p+1}\) of \(\mathrm{L}_{n+1,p+1}(\mathbf{C})\) consisting of systems of \(p+1\) vectors which form an orthonormal Hermitian basis of the vector subspace they generate; that is to say, \(\mathrm{W}_{n+1,p+1}\) consists of the matrices \(X=(x_{ij})\) which satisfy the conditions

\[
\sum_ {j = 0} ^ {n} x _ {i j} \bar {x} _ {i j} = \mathrm{I} \quad (\mathrm{I} \leqslant i \leqslant p + \mathrm{I}),
\]

\[
\sum_ {j = 0} ^ {n} x _ {i j} \overline {{{x}}} _ {k j} = 0 \quad (i \neq k).
\]

The proof of Proposition 8 of Chapter VI, § 3, no. 5 extends unaltered for the space \(\mathbf{P}_{n,p}(\mathbf{C})\) and shows that this space is connected and locally connected and that each of its points has a neighbourhood homeomorphic to \(\mathbf{C}^{(p+1)(n-p)}\) . Finally the proof of Proposition 9 of Chapter VI, § 3, no. 6 applies without change, and therefore the Grassmannian \(\mathbf{G}_{n,p}(\mathbf{C})\) is homeomorphic to \(\mathbf{P}_{n,p}(\mathbf{C})\) .

Remark. Most of the properties common to the real and complex number spaces (resp. projective spaces) are also valid for number spaces (resp. projective spaces) defined in the same way over the division ring of quaternions H; indeed, they are capable of extension to many other topological fields and division rings (cf.~Exercises 2 and 6).

\ExerciseGroupsStart
% semantic-fragments: legacy-input-seam
\relax{}
